\chapter*{Conclusioni e Sviluppi Futuri}

Le conclusioni sintetizzano il percorso svolto: dallo studio delle tecniche RAG
e dei modelli locali, alla progettazione del sistema per i3Guild, fino
all'implementazione del prototipo e alla valutazione delle sue configurazioni
principali. Il contributo centrale della tesi non consiste nell'addestramento
di un nuovo modello, ma nell'integrazione di un sistema RAG enterprise
on-premise in un contesto aziendale reale, con vincoli di privacy, risorse e
manutenibilità.

La sezione conclusiva dovrà evidenziare tre risultati. Primo, la possibilità di
usare il corpus delle Gilde come base di conoscenza interrogabile in linguaggio
naturale. Secondo, il ruolo delle ottimizzazioni non parametriche, come query
rewriting, filtri strutturati, deduplicazione e soglie dinamiche, nel rendere il
retrieval più adatto al dominio. Terzo, il valore e i limiti dell'estensione
agentica, utile per guidare l'utente nella proposta di nuove Gilde ma ancora
dipendente dalla qualità del corpus e dal comportamento dei modelli locali.

I limiti principali riguardano l'assenza di un glossario aziendale validato, la
mancanza di una ground truth estesa, la dimensione ridotta dei modelli eseguibili
localmente e la natura preliminare delle sperimentazioni su tecniche di
adattamento più invasive. Questi limiti non invalidano il prototipo, ma
definiscono con precisione il perimetro entro cui interpretarne i risultati.

Gli sviluppi futuri possono essere organizzati su tre direzioni: arricchimento
incrementale della base di conoscenza, costruzione di un dataset di valutazione
più ampio e validato con stakeholder aziendali, e sperimentazione più
approfondita di strategie come fine-tuning leggero, LoRA/QLoRA o orchestrazione
agentica più flessibile. Tali estensioni potranno essere considerate solo dopo
aver consolidato il sistema RAG di base e il processo di valutazione.
